\section{Exposing the environment and retaining movable groups}
\label{group:section}

The bridge estimates provide successful local trials.  We now select
many disjoint groups of labels and expose everything outside one group.
The two outputs are different: most cells retain many useful labels in
an unconditional probabilistic sense, while at each fixed exposed datum
the unexposed law has an exact constrained-product formula.  Keeping
these statements separate is essential for the later likelihood comparison.

\subsection{Eligible labels and their local trial laws}

The endpoints $Y,Z$ and path regularity are those of
Section~\ref{path:section}.  A label $i$ is \emph{eligible at $v$} if
\begin{enumerate}[label=\textup{(\roman*)}]
 \item $Y_i\in B_v$ and $d(Z_i,v)\le3$;
 \item $d(Y_i,Y_j)>s_0$ and $d(Z_i,Z_j)>s_0$ for every $j\ne i$;
 \item its trajectory is regular and remains in the radius-$10$ ball
 about $v$.
\end{enumerate}
The eligible sets at distinct sites are disjoint, since the starting
cells form a partition.  When $K$ is large, the confinement ball lifts
injectively to Euclidean space, so the endpoints have unique lifts
near the chosen lift of $v$.

For each eligible label $i$ at $v$, a \emph{target cell} is $B_v$ or
one of its six face-neighbor cells.  Write $Q_i^u$ for the probability
law of two independent Euclidean bridges, each of duration $t$, from
$Y_i$ to $u$ and from $u$ to $Z_i$, projected to the torus.  All lifts
in this definition are near $v$.  For a collection $F$ of fixed
obstacle paths with distinct labels other than $i$, put
\begin{equation}\label{group:success-event}
 \begin{split}
 C_i(F)=\{&\text{the trial is regular and stays within distance $10$ of $v$;}\\
          &\text{it avoids every path in $F$ at distance $r$}\}.
 \end{split}
\end{equation}
Avoidance throughout means strict distance greater than $r$ at every
time.  The endpoints are fixed in this notation.

\begin{lemma}[Useful labels in a certified cell]\label{group:certified-labels}
At a cell with a certificate there are at least $\ell D/2$ eligible
labels.  After $t$ has been chosen sufficiently small and $D$ sufficiently
large, at least nine tenths of these labels satisfy the following
stronger property: for every target cell $B_w$,
\begin{equation}\label{group:full-success}
 \left|\left\{u\in B_w:
 Q_i^u(C_i(F_i^{\rm all}))\ge0.9\right\}\right|\ge0.99.
\end{equation}
Here $F_i^{\rm all}$ consists of all other trajectories in the sampled
configuration, and $|\cdot|$ denotes spatial volume.
\end{lemma}

\begin{proof}
Start with the at least $\ell D$ labels having $Y_i\in B_v$.
Every such path visits the radius-$12$ ball and is therefore regular.
Discard the fewer than $\theta D$ visiting paths of diameter greater
than one.  The remaining paths stay within distance $2$ of $v$ and
end within distance $2$ of $v$.  Discard also those whose initial or
terminal endpoint has another endpoint within distance $s_0$.
The certificate bounds these two losses by $\theta D$ each, since
the endpoints in question lie in the radius-$5$ ball about $v$.
The remaining labels are eligible, and their number is at least
$(\ell-3\theta)D\ge\ell D/2$.

We now apply Corollary~\ref{bridge:success-supply} to the
\emph{entire} eligible list at this site, not merely the subset just
used to prove its lower cardinality bound.  Every eligible label has
its initial endpoint in $B_v$, terminal endpoint within distance
three of $v$, and both endpoints separated from all corresponding
endpoints of other labels by more than $s_0$.  The certificate supplies
at most $H_1D$ visiting obstacles, all regular.  These are precisely
the hypotheses of that corollary.  It gives
Equation~\eqref{group:full-success} for at least nine tenths of the
entire eligible list, with choices uniform in the certificate.
\end{proof}

\subsection{Disjoint groups and the information retained}

Choose an integer $M\ge1$, to be fixed after the one-path comparison
constant below, and put
\begin{equation}\label{group:number}
 G_D=\left\lfloor\frac{\ell D}{4M}\right\rfloor.
\end{equation}
At a site with at least $MG_D$ eligible labels, choose an ordered list
of $MG_D$ distinct eligible labels uniformly without replacement,
and divide it into $G_D$ successive groups of size $M$.
If there are fewer, all groups at that site are empty.  The allocation
randomizations are independent between sites conditional on the full
path sample, and use the eligible lists only.

Fix one group index $g\in\{1,\ldots,G_D\}$.  Let
$\mathcal I$ be the union of its labels over all sites.
The \emph{exposed datum} $\mathcal B$ comprises
\begin{itemize}
 \item both endpoint configurations, every eligibility indicator,
 and all allocation information;
 \item the complete trajectories of labels outside $\mathcal I$.
\end{itemize}
These last trajectories form the frozen family $F$.
The law of $\mathcal B$ is the pushforward of the stationary path law
and the allocation randomizations.  No certificate event is included
in this datum.  Subsequent expectations over $\mathcal B$ always use
this unconditional law.

For $i\in\mathcal I$ at site $v$, let $B_i$ be the unconditioned
Euclidean bridge law of duration $2t$ between its near endpoint lifts.
Define
\begin{equation}\label{group:reference}
 P_i(\dd\omega_i)=
 \frac{\ind_{C_i(F)}\,B_i(\dd\omega_i)}{a_i},
 \qquad a_i=B_i(C_i(F)).
\end{equation}
The event $C_i(F)$ includes the individual regularity and confinement
cutoffs, as in Equation~\eqref{group:success-event}.

\begin{proposition}[Exact conditional law]\label{group:conditional-law}
For almost every exposed datum the normalizers in
Equation~\eqref{group:reference} are positive, and the conditional law
of the unexposed trajectories is
\begin{equation}\label{group:product-law}
 \lambda_{\mathcal I}(\dd\omega_{\mathcal I})
 =\frac{1}{Z_{\mathcal I}}
   \ind_{\{\text{all pairs in $\mathcal I$ avoid each other}\}}
   \prod_{i\in\mathcal I}P_i(\dd\omega_i),
 \qquad Z_{\mathcal I}>0.
\end{equation}
All reference trajectories obey their own regularity and confinement
cutoffs.  Each label has at most $C M$ other labels in this group with
which a collision can be possible, where $C$ is a fixed geometric
constant.
\end{proposition}

\begin{proof}
Condition first on the endpoints.  By Lemma~\ref{path:slab}, the
free reference is a product of torus bridges and the remaining
interaction is the hard-avoidance indicator.  At fixed endpoints,
eligibility of an individual path consists of its own regularity and
confinement conditions: its starting cell, endpoint displacement and
endpoint separation are already fixed.  The probabilities of all
allocations depend only on the recorded eligibility indicators, and
are consequently constant on the conditional fiber.

Now expose all paths outside $\mathcal I$.  Every retained label was
eligible, so the only remaining individual restrictions are exactly
its two cutoffs and avoidance of the fixed frozen paths.  Confinement
to its radius-$10$ ball selects a single Euclidean lift when $K$ is
large.  The winding weights of the original torus bridge then depend
only on the fixed endpoints and cancel in normalization.  The remaining
nonfactorizing restriction is mutual avoidance inside $\mathcal I$.
This proves Equation~\eqref{group:product-law} by disintegration on
the continuous-path spaces.  Equivalently, integrate the asserted
formula against bounded functions of the exposed and unexposed
variables and use Fubini; all factors just identified reproduce
Equation~\eqref{path:law} and the allocation probabilities.

The integral of the conditional density is positive for almost every
sampled datum.  Its positivity implies positivity of each $a_i$ and
of $Z_{\mathcal I}$; otherwise that fiber would have zero weight.
Finally a retained trajectory stays within distance ten of its home
site.  Two such tubes can meet at distance $r<1$ only when their home
sites are within a fixed lattice distance.  There are at most $M$
labels per site and a fixed number of such sites, proving the last
assertion.  In particular Equation~\eqref{group:product-law} is a
conditioned product, not a claim that its trajectories are independent.
\end{proof}

\subsection{Accessibility and the probability of bad cells}

For $i\in\mathcal I$ and one of its target cells $B_w$, set
\begin{equation}\label{group:targets}
 T_i(w)=\{u\in B_w:Q_i^u(C_i(F))\ge0.9\}.
\end{equation}
A label is \emph{accessible} if $|T_i(w)|\ge0.99$ for all seven
target cells.  This property and all the sets $T_i(w)$ depend only
on $\mathcal B$.  Let $\mathcal A_v$ be the accessible labels of the
fixed group at $v$.  A site is \emph{good} if
$|\mathcal A_v|\ge M/2$, and \emph{bad} otherwise.  The sets in
Equation~\eqref{group:targets} are measurable: bridge kernels are
measurable in their endpoints, and regularity, confinement and strict
avoidance are Borel events of continuous paths.

\begin{proposition}[Unconditional supply of good cells]\label{group:supply}
For any prescribed $p>0$, the constants before $M$ can be chosen as
above, and then $M$ can be chosen sufficiently large, so that for
sufficiently large $D$ and $K$,
\begin{equation}\label{group:bad-bound}
 \PP_{\mathcal B}\{A\subset\{\text{bad sites}\}\}\le p^{|A|}
\end{equation}
for every set $A$ of distinct sites.  The bound is uniform in the
group index, the ground function, and the allowed $\alpha$-window.
It remains valid if $M$ is increased before increasing the lower
threshold on $D$.
\end{proposition}

\begin{proof}
Choose $p_1<p/2$ in Proposition~\ref{path:certificates} and meet also
the requirements of Lemma~\ref{group:certified-labels}.  Fix the
full path sample before the allocation randomizations.  At a certified
cell there are at least $\ell D/2\ge MG_D$ eligible labels, and
at least nine tenths satisfy Equation~\eqref{group:full-success}.
For a fixed group, its $M$ labels form a uniform sample without
replacement from this list.  For any $k$ specified positions in that
sample, the probability that all $k$ lack the stronger property is
at most $10^{-k}$.  A union bound therefore gives
\begin{equation}\label{group:sampling-failure}
 \PP\{\text{fewer than $M/2$ stronger labels in the group}
            \mid\text{full sample}\}
 \le q_M:=2^M10^{-M/2}.
\end{equation}
This tends to zero exponentially in $M$.

A label with the stronger property is accessible after allocation:
its frozen obstacle family is a subset of all other sampled paths,
so removing obstacles can only increase its success probabilities.
Moreover the auxiliary sampling tests at distinct certified cells
are independent conditional on the full sample.  For each subset
$A_0\subset A$, designate the cells of $A_0$ as lacking certificates
and the other cells as certified but failing the sampling test.
The probability of this event is at most
$p_1^{|A_0|}q_M^{|A\setminus A_0|}$, by conditioning first on the
full sample and then using Equation~\eqref{path:certificate-bound}.
Summing over $A_0$ bounds the probability in
Equation~\eqref{group:bad-bound} by $(p_1+q_M)^{|A|}$.
Take $M$ large enough that $q_M<p/2$.  The argument applies to each
group without conditioning on its tag or on its unexposed trajectories.
\end{proof}

\subsection{One-path comparisons at fixed exposed data}

The remaining estimates apply at a fixed datum, even when its
unexposed trajectories no longer satisfy any full-sample certificate.
Let $i$ be accessible, let $B_w$ be a target cell, and let
$U\subset T_i(w)$ be a measurable set of volume at least $0.98$.
Define the proposal $\Pi_i^U$ by choosing $u$ uniformly in $U$ and
then sampling $Q_i^u$ conditioned on $C_i(F)$.  The target set is
chosen using exposed data only; it must not depend on the trajectories
of other unexposed labels.

\begin{proposition}[Density and deterministic compatibility]
\label{group:comparison}
There is a constant $G\ge2$, fixed after $t$ and independent of $D,K,M$,
such that
\begin{equation}\label{group:density-bound}
 \Pi_i^U\ll P_i,\qquad
 0\le\frac{\dd\Pi_i^U}{\dd P_i}\le G,
 \qquad \int\frac{\dd\Pi_i^U}{\dd P_i}\dd P_i=1.
\end{equation}
There is also $\eta_D\to0$, uniform in the exposed datum, with the
following property.  Fix any regular, confined trajectory of a distinct
label in $\mathcal I$, keeping its recorded endpoints.  Under either
$P_i$ or $\Pi_i^U$, the probability of colliding with that deterministic
trajectory is at most $\eta_D$.  Its label need not be accessible.
One may take
\begin{equation}\label{group:eta}
 \eta_D=C_t\bigl(tr+rD^{3\gamma}h_0^{-1/2}\bigr)
 \le C_t' D^{-1/4}.
\end{equation}
\end{proposition}

\begin{proof}
The midpoint density $g_i$ of $B_i$ is the Euclidean Gaussian with
mean $(Y_i+Z_i)/2$ and covariance $tI$.  On every target cell it has
a lower bound $c_t>0$, uniform in the eligible endpoint data.  Indeed
those endpoints are within distance three of $v$, and every target
cell lies in a fixed bounded neighborhood of $v$; the explicit
Gaussian density gives such a bound.  Disintegrating $B_i$ at its
midpoint gives the conditional laws $Q_i^u$.  Hence
\[
 a_i=\int g_i(u)Q_i^u(C_i(F))\dd u
 \ge 0.99\cdot0.9\,c_t=:a_t>0.
\]
At a successful trajectory with midpoint $u$, the proposal density
relative to $P_i$ is
\[
 \frac{a_i\ind_U(u)}{|U|g_i(u)Q_i^u(C_i(F))}
 \le \frac1{0.98\cdot0.9\,c_t},
\]
since $a_i\le1$.  Taking $G$ to be the maximum of two and this
constant proves Equation~\eqref{group:density-bound}; its final
identity holds because both measures are probabilities.

For collision with the fixed opponent, first use the unconditioned
bridge $B_i$.  The event contributing under $P_i$ satisfies the
regularity and confinement cutoffs.  Proposition~\ref{bridge:single-hit}
bounds its probability by
$C(tr+rD^{3\gamma}h_0^{-1/2})$.  Dividing by $a_i\ge a_t$
proves the required estimate for $P_i$.
For the proposal, integrate that proposition's two-leg estimate over the
whole target cell; restricting the integral to $U$ can only decrease
it.  Division by $|U|\ge0.98$ and by the through-$u$ success
probability, which is at least $0.9$, proves the proposal estimate.
The endpoint separation required in that proposition holds because $i$ was
eligible and the opponent has a distinct recorded label.

Finally $\gamma=1/12$, $h_0=D^{-1}$ and $r=\alpha/D$ give
$rD^{3\gamma}h_0^{-1/2}=\alpha D^{-1/4}$ and $tr=O_t(D^{-1})$.
The $\alpha$-window is fixed.  This proves
Equation~\eqref{group:eta}, with constants depending on the already
fixed $t$ but not on the later choice of $M$ or the dilution threshold.
\end{proof}

Each changed path has only $CM$ potential opponents by confinement.
Consequently a union of a fixed number of comparison copies has a
one-variable failure bound of the form $CM\eta_D$.  The comparison
argument will choose $M$ after $G$, then increase $D$ so that this
quantity is small.  No success probability is multiplied along the
entire length of a lattice path: the normalizers for separated changes
will cancel in the second-moment calculation.
